\section{Problem Formulation}
\label{sec:problem}

\subsection{Dense-Intelligence Tasks}

The public Argus formulation defines a \emph{dense-intelligence task} as one that
sustains high-frequency reasoning, tool use, verification, and iteration across a
continuous time window until it produces a measurable result~\citep{argus_site}.
Three conditions determine whether a task has this shape:

\begin{enumerate}[label=\textbf{T\arabic*.}]
  \item \textbf{Invention.} The answer is not directly retrievable from a manual or
    database; it must be discovered through search and feedback.
  \item \textbf{Long horizon.} One pass is insufficient. Progress requires many
    dependent iterations over hours or days.
  \item \textbf{Verifiability.} A task-native evaluator can distinguish genuine
    improvement from a persuasive but incorrect claim.
\end{enumerate}

The defining loop is therefore proposal $\rightarrow$ execution $\rightarrow$
measurement $\rightarrow$ revised proposal. Performance engineering, security
research, scientific search, and quantitative research differ in domain but share
this operational structure. The task definition excludes both open-ended ideation
without an evaluator and fixed workflows that require time but no new decisions.

We model a campaign by an objective $o$, an initial artifact $y_0$, an evaluation
procedure $f$, and a resource budget $B$. The artifact may be code, a model, a
dataset, an experimental harness, a proof, or a manuscript. The campaign is divided
into bounded missions indexed by $t$, each producing

\begin{equation}
  \tau_t=(s_t,a_t,y_t,e_t,r_t),
  \label{eq:trajectory}
\end{equation}

where $s_t$ is the mission state, $a_t$ the actions taken, $y_t$ the resulting
artifact state, $e_t$ the observed measurements, and $r_t$ the review outcome. A mission
must terminate, pause, or transfer control explicitly; continuity belongs to the
persistent campaign state rather than an unbounded model session.

\subsection{Dense-Intelligence Density}

The website expresses useful intelligence per unit wall-clock time as

\begin{equation}
  \rho_I(T)
  =\frac{1}{T}\int_0^T
  \dot N_{\mathrm{tok}}(t)\,
  \eta_r(t)\,\eta_a(t)\,\eta_v(t)\,dt.
  \label{eq:dense-density}
\end{equation}

$\dot N_{\mathrm{tok}}(t)$ is the instantaneous rate of token-mediated reasoning
and action. The factors $\eta_r(t)$, $\eta_a(t)$, and $\eta_v(t)$ represent the
fractions converted into relevant reasoning, effective action, and valid
verification. The product is intentionally stricter than token throughput: a run
that repeatedly reads the same state, emits activity without changing an artifact,
or accepts unchecked claims can consume many tokens while contributing little to
$\rho_I$.

The public page summarizes the intended continuous-utilization advantage over a
fixed 48-hour window as

\begin{equation}
  \rho_I^{\mathrm{Argus}}(48\mathrm{h})
  \gg
  \rho_I^{\mathrm{human}}(48\mathrm{h}).
  \label{eq:density-hypothesis}
\end{equation}

Equation~\ref{eq:density-hypothesis} is a public design hypothesis, not a measured
result in this report. It refers to the density of sustained work over continuous
time, not to unconditional superiority of any individual reasoning step. Neither
$\rho_I$ nor its conversion factors are instrumented as a benchmark score in the
current evaluation.

\subsection{Online Life and Runtime Evolution}

The second website object separates model parameters from the operating state
around the model:

\begin{equation}
  \begin{aligned}
    H_{t+1} &= U(H_t,\tau_t,E_t), \\
    \theta_{t+1} &= \theta_t.
  \end{aligned}
  \label{eq:runtime-update}
\end{equation}

with the public state definition

\begin{equation}
  H_t=\left\{\text{Memory},\text{Skills},\text{Tools},
  \text{Verifiers},\text{Routing}\right\}.
  \label{eq:website-state}
\end{equation}

$E_t$ is the subset of results admitted from mission $t$. The equality for
$\theta$ is a scope condition: online evolution does not require a gradient update
to the underlying model. Section~\ref{sec:method} refines $H_t$ into named ownership
surfaces and tracks task/evaluation definitions as an additional component $Q_t$.
The update $U$ is partial; many missions change only memory, and some change no
reusable component.

\subsection{Capability-Set Expansion}

Let $C_t$ denote the effective capability set available at time $t$. The website
writes review-gated admission and the desired breadth/depth behavior as

\begin{equation}
  \begin{aligned}
    C_{t+1}
      &= C_t\cup\{c:\operatorname{Verify}(c,E_t)\geq\epsilon\}, \\
    W_{t+1} &\geq W_t,
    & D_{t+1}(d) &\geq D_t(d).
  \end{aligned}
  \label{eq:ood-expansion}
\end{equation}

where $W_t$ is frontier width across domains and $D_t(d)$ is depth within domain
$d$. The inequalities are retention desiderata over validated capability, not an
empirical guarantee that every run improves the system. Skills can become stale,
verifiers can be revised, and performance can regress on harder instances. A
negative result can still add value by excluding a failed branch even when the
effective capability set does not grow.

\subsection{Research Value}

Dense activity matters only when it produces new, valid, reusable information. The
website therefore defines the accumulated research value over horizon $T$ as

\begin{equation}
  V_R(T)
  =\int_0^T
  \rho_I(t)\,\Delta I(t)\,p_{\mathrm{valid}}(t)\,
  \eta_{\mathrm{reuse}}(t)\,dt,
  \label{eq:research-value}
\end{equation}

where $\Delta I(t)$ is information gain, $p_{\mathrm{valid}}(t)$ is the probability
that the claimed increment survives review, and
$\eta_{\mathrm{reuse}}(t)$ is the fraction that remains useful for later work.
Equation~\ref{eq:research-value} explains why token volume alone is not the product:
research value collapses if the work is repetitive, invalid, or non-reusable. It is
a conceptual objective, not a calibrated monetary or benchmark metric in this
report.

\subsection{The Information Advantage of Process Data}

Finally, the public formalization distinguishes the accepted artifact from the
trajectory that produced it:

\begin{equation}
  D_{\mathrm{process}}
  =\{(s_k,a_k,e_k,r_k,\Delta H_k)\}_{k=1}^{N}
  \supsetneq
  D_{\mathrm{final}}=\{y^\star\}.
  \label{eq:process-data}
\end{equation}

The process record contains states, actions, measurements, review feedback, and
runtime-state changes, including failed branches omitted from $y^\star$. This
additional information can reduce repeated search, support later analysis, and provide
grounded material for later skill construction or evaluation. \system{} retains
typed, credential-redacted observations and verdicts, not private chain-of-thought
transcripts.
